\documentclass[12pt, a4paper]{article}

% --- Packages ---
\usepackage[margin=2.5cm]{geometry}
\setlength{\headheight}{28.44pt}
\addtolength{\topmargin}{-16.44pt}
\usepackage{setspace}
\usepackage{booktabs}
\usepackage{array}
\usepackage{longtable}
\usepackage{tabularx}
\usepackage{hyperref}
\usepackage[dvipsnames,table]{xcolor}
\usepackage{titlesec}
\usepackage{fancyhdr}
\usepackage{graphicx}
\usepackage{amsmath}
\usepackage{amssymb}
\usepackage{enumitem}
\usepackage{parskip}
\usepackage{microtype}
\usepackage[T1]{fontenc}
\usepackage{lmodern}
\usepackage[polish]{babel}
\usepackage[utf8]{inputenc}

\definecolor{lightgray}{gray}{0.92}

\addto\captionspolish{
    \renewcommand{\figurename}{Rycina}
}

% --- Hyperref setup ---
\hypersetup{
    colorlinks = true,
    linkcolor  = blue!60!black,
    citecolor  = blue!60!black,
    urlcolor   = blue!60!black
}

% --- Header/Footer ---
\pagestyle{fancy}
\fancyhf{}
\rhead{\textit{Notatka Przesiewowa -- Analiza Związków}}
\lhead{
\includegraphics[scale=0.2]{graphics/logo-i-orzel.png}}
\rfoot{Strona \thepage}
\lfoot{\textit{}}

% --- Section formatting ---
\titleformat{\section}{\large\bfseries}{{\thesection}}{1em}{}[\titlerule]
\titleformat{\subsection}{\normalsize\bfseries}{{\thesubsection}}{1em}{}

% --- Custom commands ---
\newcommand{\placeholder}[1]{\textcolor{red!70!black}{[#1]}}
\newcommand{\TODO}[1]{\textcolor{orange!80!black}{\textbf{TODO:} #1}}

% ============================================================
\begin{document}
% ============================================================

% --- Strona tytułowa ---
\hypersetup{pageanchor=false}
\begin{titlepage}
    \centering
    \vspace*{2cm}
    {\LARGE \textbf{Notatka Przesiewowa} \par}
    \vspace{0.5cm}
    {\large Wstępna identyfikacja szkodliwych związków chemicznych\\
    w badaniu przeżycia róż\par}
    \vspace{2cm}
    
\includegraphics[width=0.3\textwidth]{graphics/UMB-logo.png}\par\vspace{2cm}

    \begin{tabular}{ll}
        \toprule
        \textbf{Wersja dokumentu}  & 1.0 \\
        \textbf{Data}              & 29 maja 2026 \\
        \textbf{Autorzy}           & Julia Czyż, Urszula Skuńczyk, Paulina Włodarczyk \\
        \textbf{Afiliacja}         & Uniwersytet Medyczny w Białymstoku \\
                                   & Zakład Biostatystyki i Informatyki Medycznej \\
        \bottomrule
    \end{tabular}

    \vspace{2cm}
    \vfill
\end{titlepage}

\hypersetup{pageanchor=true}

% --- Spis treści ---
\tableofcontents
\newpage

\onehalfspacing

% ============================================================
\section{Cel i kontekst}
% ============================================================
Niniejsza notatka stanowi dokument poprzedzający analizę pilotażową oraz analizę główną badania dotyczącego skuteczności związków chemicznych w konserwacji kwiatów. Jej celem jest identyfikacja związków o potencjalnie szkodliwym działaniu oraz formalne uzasadnienie decyzji o wykluczeniu ich z właściwego badania i analiz z nim związanych. 

Przeprowadzenie owego badania przesiewowego (ang. \textit{screening}) ma na celu optymalizację wykorzystania zasobów oraz ograniczenie niepotrzebnych strat biologicznych
-- róże nie będą przydzielane do związków wykazujących szkodliwe działanie. Pozwoli to uniknąć zbędnej śmiertelności roślin w analizie pilotażowej i analizie głównej.

% ============================================================
\section{Badane związki}
% ============================================================

W badaniu uwzględniono łącznie 15 związków chemicznych. Jako punkt odniesienia (kontrola)
przyjęto wodę destylowaną (związek 1), wyróżnioną kolorem szarym. Tabela~\ref{tab:zwiazki} zawiera zestawienie wszystkich
badanych substancji.

\vspace{0.5em}
\begin{table}[ht]
\centering
\caption{Zestawienie badanych związków chemicznych}
\label{tab:zwiazki}
\renewcommand{\arraystretch}{1.3}
\begin{tabularx}{\textwidth}{>{\centering\bfseries}p{1.5cm} >{\centering\arraybackslash}X}
\toprule
\textbf{ID} & \textbf{Nazwa związku} \\
\midrule
\rowcolor{lightgray}
1  & Distilled Water           \\
2  & Apathic Acid              \\
3  & Beerse Brew               \\
4  & Concentrate of Caducues   \\
5  & Distillate of Discovery   \\
6  & Essence of Epiphaneia     \\
7  & Four in December          \\
8  & Granules of Geheref       \\
9  & Kar-Hamel Mooh            \\
10 & Lucifer's Liquid          \\
11 & Noospherol                \\
12 & Oil of John's Son         \\
13 & Power of Perlimpinpin     \\
14 & Spirit of Scienza         \\
15 & Zest of Zen               \\
\bottomrule
\end{tabularx}
\end{table}

% ============================================================
\section{Dane}
% ============================================================

Analiza przesiewowa przeprowadzona została na danych pilotażowych obejmujących
$n = 900$ obserwacji (60 kwiatów na związek, $n_c$ = 60). Dla każdej jednostki eksperymentalnej zarejestrowano:

\begin{itemize}[noitemsep]
    \item \texttt{compound} -- identyfikator zastosowanego związku (1--15),
    \item \texttt{time} -- czas przeżycia kwiatu od moment zanurzenia w roztowrze do dnia utraty walorów estetycznych (w dniach),
    \item \texttt{can.be.presented} -- wskaźnik zdarzenia: \texttt{FALSE} (1) oznacza
          wystąpienie zdarzenia (obumarcie kwiatu), \texttt{TRUE} (0) -- obserwację prawostronnie cenzurowaną.
\end{itemize}

W danych pilotażowych odsetek obserwacji cenzurowanych wyniósł \textbf{0\%}, co oznacza, że wszystkie 900 kwiatów obumarło w trakcie 30 dni eksperymentu.
Brak cenzurowania oznacza, że test log-rank dysponuje pełną informacją o czasach
przeżycia wszystkich jednostek ($n_c$ = 60), co przekłada się na wysoką moc testu.

% ============================================================
\section{Metody}\label{sec:metody}
% ============================================================
Dla każdego związku $c \in \{2, 3, \ldots, 15\}$
przeprowadzany jest jednostronny, nieparametryczny test log-rank względem grupy kontrolnej ($c = 1$),
z następującymi hipotezami:

\begin{align*}
    H_0 &: S_c(t) = S_1(t) \quad \text{dla wszystkich } t \geq 0 \\
    H_1 &: S_c(t) < S_1(t) \quad \text{dla pewnego } t > 0
\end{align*}

gdzie $S_c(t)$ oznacza funkcję przeżycia dla związku $c$, a $S_1(t)$ -- funkcję przeżycia
grupy kontrolnej (woda destylowana). Hipoteza zerowa $H_0$ odnosi się do sytuacji, w której obserwuje się brak istotnych różnic między krzywymi przeżycia dla obu badanych substancji, natomiast hipoteza alternatywna $H_1$, że zastosowanie
związku $c$ skraca czas przeżycia kwiatu w porównaniu z kontrolą.

W celu kontroli współczynnika FWER (ang. \textit{Family-Wise Error Rate}) stosowana jest poprawka Holma na wielokrotne testowanie (14 równoległych porównań).

Związek uznaje się za szkodliwy, jeśli spełniony jest następujący warunek:

\begin{enumerate}[label=\arabic*)]
    \item Jednostronny test log-rank daje wartość $p < 0{,}05$ po korekcie na wielokrotne testowanie.
\end{enumerate}

Ponieważ implementacja algorytmu dostarcza dwustronnej wartości $p$, aby otrzymać jednostronną wartość $p$ dokonuje się następującego przekształcenia:
\[
    p_{\text{jednostronna}} = \begin{cases}
        p_{\text{dwustronna}} / 2 & \text{jeśli } O_c > E_c \quad \text{(związek szkodliwy } H_1\text{)} \\
        1 - p_{\text{dwustronna}} / 2 & \text{jeśli } O_c \leq E_c \quad \text{(związek neutralny lub korzystny)}
    \end{cases}
\]
gdzie $O_c$ i $E_c$ to odpowiednio liczba obserwowanych i oczekiwanych zdarzeń dla związku $c$.
W przypadku związku neutralnego lub korzystnego wartość $p_{\text{jednostr}} > 0{,}5$, co zapobiega ich błędnemu wykluczeniu.

% ============================================================
\section{Analiza statystyczna}
% ============================================================

\subsection{Krzywe Kaplana--Meiera}

Na Rycinie~\ref{fig:km} przedstawiono krzywe przeżycia Kaplana--Meiera dla wszystkich
15 związków. W celu zwiększenia przejrzystości oraz ułatwienia oceny działania poszczególnych związków od 2 do 14 zestawiono je na wykresach z grupą kontrolną -- wodą destylowaną, oznaczoną szarą linią przerywaną. Krzywe przeżycia związków 2 (\textit{Apathic Acid}), 3 (\textit{Beerse Brew}) oraz 10 (\textit{Lucifer's Liquid}) przebiegają wyraźnie poniżej krzywej grupy kontrolnej, co wskazuje na znacznie gorsze właściwości konserwacji kwiatów w porównaniu z wodą destylowaną.

\begin{figure}[!h]
    \centering
    \includegraphics[width=0.9\textwidth]{plots/km_compounds_facet.pdf}
    \caption{Krzywe przeżycia Kaplana--Meiera dla 15 związków chemicznych i grupy kontrolnej.}
    \label{fig:km}
\end{figure}
\newpage
\subsection{Wyniki testów log-rank}

Tabela~\ref{tab:logrank} zawiera mediany czasu przeżycia oraz wyniki jednostronnych testów log-rank dla każdego z 14 porównań z grupą kontrolną wraz z wartościami $p$ po zastosowaniu korekty Holma na wielokrotne testowanie. Związki uznane za szkodliwe wyróżniono kolorem czerwonym.

\vspace{0.5em}
\begin{table}[ht]
\centering
\caption{Wyniki analizy przesiewowej -- jednostronne testy log-rank vs kontrola.}
\label{tab:logrank}
\renewcommand{\arraystretch}{1.3}
\begin{tabularx}{\textwidth}{>{\centering}p{1.5cm} >{\centering}p{3.5cm} >{\centering}X >{\centering\arraybackslash}X}
\toprule
\textbf{ID} & \textbf{Mediana KM (dni)} & \textbf{$p_{\text{jednostr}}$} & \textbf{$p_{\text{Holm}}$} \\
\midrule
1  & 9{,}0  & --          & -- \\
\rowcolor{red!20}
2  & 3{,}0  & $< 0{,}001$ & $< 0{,}001$ \\
\rowcolor{red!20}
3  & 3{,}0  & $< 0{,}001$ & $< 0{,}001$ \\
4  & 17{,}5 & $> 0{,}999$ & $> 0{,}999$ \\
5  & 20{,}0 & $> 0{,}999$ & $> 0{,}999$ \\
6  & 17{,}0 & $> 0{,}999$ & $> 0{,}999$ \\
7  & 13{,}0 & $> 0{,}999$ & $> 0{,}999$ \\
8  & 20{,}0 & $> 0{,}999$ & $> 0{,}999$ \\
9  & 15{,}0 & $> 0{,}999$ & $> 0{,}999$ \\
\rowcolor{red!20}
10 & 1{,}0  & $< 0{,}001$ & $< 0{,}001$ \\
11 & 21{,}0 & $> 0{,}999$ & $> 0{,}999$ \\
12 & 19{,}5 & $> 0{,}999$ & $> 0{,}999$ \\
13 & 20{,}0 & $> 0{,}999$ & $> 0{,}999$ \\
14 & 14{,}5 & $> 0{,}999$ & $> 0{,}999$ \\
15 & 17{,}0 & $> 0{,}999$ & $> 0{,}999$ \\
\bottomrule
\end{tabularx}
\end{table}

% ============================================================
\section{Podsumowanie}
% ============================================================

Na podstawie analizy przeprowadzonej zgodnie z kryterium określonym w sekcji~\ref{sec:metody}, na przyjętym poziomie istotności $\alpha = 0{,}05$, uznano następujące związki za szkodliwe:

\begin{itemize}
    \item 2 -- Apathic Acid,
    \item 3 -- Beerse Brew,
    \item 10 -- Lucifer's Liquid.
\end{itemize}

Związki te zostają \textbf{wykluczone ze wszystkich dalszych etapów badania}, tj.\
z analizy pilotażowej oraz analizy głównej. Pozostałe 12 związków (w tym grupa kontrolna) zostają przekazane do analizy pilotażowej zgodnie z protokołem.

\end{document}
